%% kapa: exam=1 round=27 year=2016 period=2 form=A subject=appraisal_law subject_ko=감정평가관계법규 seq=28 number=28
%% source: Q-Net 감정평가사 기출문제 (cst00309) artlSeq=5169402 file=extracted/1008/5169402/제27회 감정평가사 1차 문제/제27회 감정평가사 2교시문제(A형).pdf
%% conversion: kordoc 4.24.51 -> tools/gen1.py (auto); review=proofread-sonnet; flags=none
\begin{question}{28}
공간정보의 구축 및 관리 등에 관한 법령상 지번의 구성과 부여방법에 관한 설명으로 옳지 않은 것은?
\begin{choices}
  \item 지번은 북서에서 남동으로 순차적으로 부여하여야 한다.
  \item 토지소유자가 합병 전의 필지에 주거ㆍ사무실 등의 건축물이 있어서 그 건축물이 위치한 지번을 합병 후의 지번으로 신청한 경우에도 합병 대상 지번 중 선순위의 지번으로 부여하여야 한다.
  \item 분할의 경우에는 분할 후의 필지 중 주거․사무실 등의 건축물이 있는 필지에 대해서는 분할 전의 지번을 우선하여 부여하여야 한다.
  \item 지번은 아라비아숫자로 표기하되, 임야대장 및 임야도에 등록하는 토지의 지번은 숫자 앞에 "산"자를 붙인다.
  \item 신규등록 및 등록전환의 경우에 대상토지가 여러 필지로 되어 있는 경우에는 그 지번부여지역의 최종 본번의 다음 순번부터 본번으로 하여 순차적으로 지번을 부여할 수 있다.
\end{choices}
\end{question}
\answer{2}
